\chapter{Introdução} \label{ch:introduccion}

% ============================================================
%  BLOCO A — RELEVÂNCIA DO FENÔMENO
%  Propósito: Estabelecer por que os ciclones extratropicais
%  do Atlântico Sul importam em escala global e regional.
%  Justifica a escolha do objeto de estudo antes de
%  qualquer revisão bibliográfica.
% ============================================================

% --- A1: Papel climático e dinâmico global ---
Sob uma perspectiva dinâmica global, a análise dos \emph{storm tracks} realizada por \citet{hoskins2005new} mostra os ciclones extratropicais como o mecanismo dominante para o transporte latitudinal de energia no Hemisfério Sul, equilibrando o gradiente térmico planetário. A quantificação desses transportes foi possível graças ao desenvolvimento de algoritmos de rastreamento objetivo, em que \citet{sinclair1994objective} estabeleceu o precedente fundamental ao utilizar a vorticidade relativa para identificar e rastrear sistemas independentemente do escoamento de fundo.

% --- A2: Impacto regional e relevância societal ---
No hemisfério sul, esses sistemas representam o fenômeno de severidade geofísica dominante para a zona costeira regional \citep{gramcianinov2023impact}. No litoral sul do Brasil, essa ameaça se traduz em riscos meteoceanográficos severos e perdas humanas documentadas: \citet{bartolomei2024extremos} reportam a letalidade causada por ciclones extratropicais invernais nessa região, enquanto \citet{miranda2026patterns} confirmam que esses sistemas são os principais motores de tempestades costeiras no Atlântico Sul. Sua influência se estende ainda ao estado do mar: \citet{sasaki2021intraseasonal} demonstram que esses ciclones modulam a variabilidade intrassazonal das ondas no Atlântico Sul ocidental, onde a configuração do vento em superfície determina a transferência de energia para o estado de mar, amplificando o risco costeiro além da duração do evento sinótico.

% ============================================================
%  BLOCO B — ESTADO DA ARTE: O QUE SE SABE
%  Propósito: Sintetizar o conhecimento consolidado em quatro
%  dimensões: (B1) climatologia de trajetórias, (B2)
%  diversidade estrutural, (B3) forçante orográfico e
%  regionalização, e (B4) papel da umidade e dos
%  transportes diabáticos. Este bloco constrói o arcabouço
%  conceitual necessário para delimitar a lacuna no Bloco C.
% ============================================================

% --- B1: Climatologia de trajetórias e ciclogênese ---
No contexto regional, \citet{reboita2010south} forneceram uma das primeiras climatologias para o Atlântico Sul, identificando três centros preferenciais de ciclogênese vinculados à interação entre o escoamento perturbado pelos Andes e os gradientes térmicos oceânicos. Investigações posteriores como \citet{gramcianinov2020analysis} e pesquisas recentes como \citet{padilhareinke2026characterization} caracterizaram com robustez a distribuição espaço-temporal desses sistemas, consolidando uma base climatológica sobre a frequência e localização da atividade ciclônica na região.

% --- B2: Diversidade estrutural (tipologia e ciclo de vida) ---
Não obstante, a natureza dessas perturbações na região é heterogênea. \citet{evans2012climatology} destacam a prevalência de ciclones subtropicais com estruturas térmicas híbridas, e \citet{hart2003cyclone} demonstrou que os sistemas transitam ao longo de um contínuo estrutural em vez de pertencerem a categorias estáticas. De forma complementar, \citet{gozzo2013air} documentam a ocorrência de sistemas que seguem o modelo estrutural de Shapiro-Keyser, caracterizados por uma marcada fratura da frente fria e o desenvolvimento de uma seclusão quente durante sua etapa de maturidade. Essas variações morfológicas, respaldadas por climatologias regionais de longo prazo \citep{gozzo2014subtropical}, evidenciam que o Atlântico Sul se afasta frequentemente da conceituação de ciclones convencionais.

% --- B3: Forçante orográfico e regionalização ---
Compreender a gênese desses sistemas exige considerar o papel do continente sul-americano como perturbação orográfica. A Cordilheira dos Andes interage mecanicamente com o escoamento dos ventos de oeste, forçando a compressão da coluna de vorticidade a barlavento e seu estiramento vertical a sotavento, um mecanismo físico descrito por \citet{gan1994influence} que favorece a formação de centros de baixa pressão a leste da cordilheira. Esse forçante topográfico, combinado com a baroclinicidade costeira e os gradientes de temperatura superficial do mar, consolida focos de atividade ciclogenética na costa leste da América do Sul \citep{reboita2026meteorology}.

\label{sec:definicion_regiones}%
Adotando a regionalização proposta por \citet{gramcianinov2019properties}, este estudo se concentra em três regiões-chave: a região Sul-Sudeste do Brasil (\textit{South Brazil}, SBR), a bacia de drenagem do Rio da Prata (\textit{La Plata Basin}, LPB, em inglês) e a costa sudeste da Argentina (\textit{Argentina}, ARG). Essas regiões de ciclogênese foram amplamente documentadas por diversos autores como \citet{reboita2010south, crespo2020potential, gramcianinov2019properties, reboita2026meteorology}: (1) a região SBR, favorecida pelo efeito de sotavento dos Andes e pelo jato de baixos níveis (SALLJ), onde \citet{reboita2022from} documentam frequentes sistemas híbridos e transições tropicais; (2) a região LPB, caracterizada por ciclogêneses explosivas e influenciada pelo transporte de umidade do SALLJ, sendo determinante para a desestabilização da atmosfera \citep{crespo2020potential}; e (3) a região ARG, dominada pela baroclinicidade da frente polar, onde a dinâmica do jato polar e a propagação de ondas de Rossby constituem os principais mecanismos de desenvolvimento \citep{simmonds2000variability}. Essa estratificação responde a regimes dinâmicos distintos: enquanto os sistemas em SBR e LPB exibem uma sensibilidade projetada para a intensificação dos fluxos de umidade e a advecção quente em camadas baixas, as tendências na estrutura média dos controladores físicos para a região ARG resultam relativamente fracas \citep{gramcianinov2024early}.

% --- B4: Ingredientes dinâmicos: umidade, SALLJ e diabática ---
A variabilidade da atividade ciclônica é modulada pela corrente de jato e sua interação com as camadas baixas. Ao cruzar os Andes, segundo \citet{vera2002cold}, o jato se realinha no lado de sotavento, gerando convergência de ventos em níveis baixos e divergência em altitude, o que intensifica os ciclones e provoca precipitações intensas no sudeste da América do Sul. Simultaneamente, o transporte meridional de umidade dos trópicos é mediado pelo Jato de Baixos Níveis da América do Sul (SALLJ). Como apontam \citet{reboita2022from}, a advecção de ar quente e úmido no flanco leste do ciclone não apenas desestabiliza a atmosfera, mas atua como uma fonte diabática que, ao liberar calor latente, reforça a intensidade do sistema.
No Hemisfério Norte (HN), o escoamento ascendente da Esteira Transportadora Quente (WCB, em inglês) é descrito tipicamente em direção ao nordeste \citep{blanchard2021warm}; no Hemisfério Sul (HS), esse escoamento se orienta preferencialmente para o sudeste do sistema.
% de modo que os padrões estruturais do HN devem ser interpretados como uma imagem especular latitudinal ao serem transpostos para o HS.
A WCB não apenas produz precipitação de grande escala, mas frequentemente abriga núcleos de convecção intensa embebida \citep{oertel2021warm}, contribuindo para a produção de precipitação extrema desacoplada do centro do ciclone \citep{naud2020evaluation}. Pesquisas recentes de \citet{gentile2025response} antecipam, mediante modelagens de alta resolução, que o setor quente dos ciclones extratropicais mais intensos se intensificará, incrementando significativamente os ventos e a precipitação extrema durante suas etapas de desenvolvimento.

% ============================================================
%  BLOCO C — LACUNA DE CONHECIMENTO: O QUE NÃO SE SABE
%  Propósito: Articular de forma unificada e explícita a
%  lacuna científica que motiva esta pesquisa. A lacuna
%  se constrói em três dimensões acumulativas:
%  (C1) insuficiência das métricas escalares de intensidade,
%  (C2) ausência de caracterização estrutural por fase
%  evolutiva, e (C3) limitações de resolução dos dados.
%  Cada dimensão prepara a justificativa do Bloco D.
% ============================================================

% --- C1: Insuficiência das métricas escalares de intensidade ---
Embora existam múltiplos estudos sobre esses sistemas, persiste uma limitação conceitual na forma como a severidade desses sistemas é quantificada. Como argumenta \citet{corner2025classification} para o Atlântico Norte, a intensidade de um ciclone não pode ser capturada por uma única métrica, pois a relação entre a intensidade do sistema e os ventos em superfície não é linear. \citet{sinclair2023relationship} advertem que um ciclone pode ser dinamicamente intenso mas gerar pouca chuva se carecer do suprimento adequado de umidade, sublinhando que os máximos atingidos não necessariamente coincidem com o mínimo de pressão ocorrido. \citet{bitencourt2010relating} evidenciam que os máximos de vento não se distribuem uniformemente ao redor do centro, e \citet{cardoso2022synoptic} o confirma ao mostrar que se concentram em quadrantes específicos. No HN, \citet{eisenstein2023identification} destacam ainda que os maiores impactos em superfície costumam estar associados a estruturas de mesoescala específicas, como as esteiras transportadoras frias (CCB), cuja localização espacial em relação ao centro de vorticidade varia durante o ciclo de vida do sistema.

% --- C2: Ausência de caracterização estrutural por fase evolutiva ---
Existe uma lacuna na compreensão de como se estrutura espacialmente a distribuição de vento e precipitação durante a evolução do ciclone, e de como ela varia nas distintas fases do ciclo de vida.
Na América do Sul, trabalhos de \citet{gozzo2013air} e \citet{reboita2022from} revelaram que os ciclones exibem evoluções que se afastam dos modelos tradicionais; da mesma forma, \citet{marrafon2022classificacao} aponta que essas evoluções estão vinculadas a condições de tempo adverso.
A ocorrência espacial dos máximos de vento e precipitação dentro de uma mesma fase do ciclone é uma consequência estrutural do acoplamento dinâmico entre os mecanismos que governam cada campo \citep{chen2025characteristics}, de modo que sua análise em conjunto é um insumo necessário para caracterizar a morfologia superficial do sistema.
No entanto, seu estudo requer uma descrição da organização espacial independente de cada campo e de sua covariabilidade estrutural ao longo do ciclo de vida. Recentemente, \citet{han2025system} enfatiza que são necessários marcos de classificação que vinculem a evolução do sistema às suas manifestações extremas de vento e precipitação.

% --- C3: Limitações de resolução dos dados ---
Essa caracterização estrutural tem sido limitada pela resolução espacial dos dados disponíveis. \citet{priestley2022improved} demonstram que os modelos subestimam os ventos extremos associados a ciclones ao não resolverem explicitamente suas estruturas de mesoescala. Da mesma forma, \citet{neu2013intercomparison} apontam que a variabilidade na forma e no tamanho dos ciclones extratropicais gera incertezas na determinação de sua intensidade em superfície, especialmente em regiões com topografia complexa.

% ============================================================
%  BLOCO D — JUSTIFICATIVA E OPORTUNIDADE METODOLÓGICA
%  Propósito: Demonstrar que a lacuna identificada no Bloco C
%  é agora tratável graças à convergência de três
%  avanços recentes: (D1) a reanálise ERA5, (D2) o
%  algoritmo de classificação de fases CycloPhaser, e
%  (D3) o rastreamento por vorticidade relativa. Este bloco
%  transforma o problema em uma pesquisa viável.
% ============================================================

Este estudo utiliza a base de dados de trajetórias de \citet{coutodesouza2024new}, fundamentada no rastreamento por vorticidade relativa a 850 hPa ($\zeta_{850}$), variável que permite isolar a circulação do sistema do escoamento de fundo e capturar sistemas em etapas iniciais que os algoritmos baseados em pressão mínima tendem a omitir \citep{sinclair1994objective,sinclair1997objective}.
Os fundamentos dinâmicos e a validação desse critério para o Hemisfério Sul são desenvolvidos na Seção~\ref{sec:tb_vorticidad}. O uso de $\zeta_{850}$ é particularmente crítico para esta pesquisa porque garante a representação da fase incipiente, condição necessária para uma caracterização estrutural completa do ciclo de vida.

Para caracterizar a evolução física da estrutura interna, adota-se a segmentação objetiva desenvolvida por \citet{coutodesouza2024new} mediante o algoritmo \textit{CycloPhaser}, que define quatro fases dinâmicas discretas: incipiente, intensificação, maturidade e decaimento \citep{deSouza2025CycloPhaser}.
Esse marco supera as limitações das classificações clássicas no contexto do Atlântico Sul (discutidas na Seção~\ref{subsec:tb_fases}), permitindo analisar como se transforma a distribuição espacial dos campos de superfície ao longo das distintas etapas do sistema. Essa estratégia se alinha com propostas globais recentes como o \textit{SyCLoPS} de \citet{han2025system}.
A caracterização dessa evolução estrutural se apoia ainda no modelo de Fratura Frontal \citep{shapiro1990life,shapiro1999bridge} como marco de referência para interpretar a distribuição espacial dos máximos de vento e precipitação; seu desenvolvimento é apresentado na Seção~\ref{subsec:tb_modelos}.

A viabilidade dessa caracterização estrutural reside na disponibilidade da reanálise ERA5 \citep{hersbach2020era5}, cuja resolução espacial ($\sim$31 km) e temporal (horária) permite resolver os gradientes de pressão e os fluxos de umidade que gerações anteriores de dados suavizavam. Como demonstram \citet{priestley2022improved} e \citet{dalanhese2023new}, essa maior resolução reduz a subestimação da intensidade dos ciclones observada em produtos de gerações anteriores e permite capturar a organização espacial dos gradientes de vento e os núcleos de precipitação. Validações específicas para o Atlântico Sul realizadas por \citet{gramcianinov2020analysis} confirmam que o ERA5 oferece uma representação superior da intensidade dos ventos superficiais e da estrutura dos ciclones em comparação com esses produtos anteriores. No entanto, o ERA5 apresenta vieses documentados na representação dos extremos mais intensos, cuja magnitude e consequências para a interpretação dos resultados são discutidas na Seção~\ref{sec:datos_era5}.

% ============================================================
%  BLOCO E — PERGUNTA DE PESQUISA, HIPÓTESE E
%             ESTRATÉGIA ANALÍTICA
%  Propósito: Formular o problema central e a hipótese
%  de trabalho, e introduzir imediatamente a estratégia
%  estatística (MEOF) que torna a hipótese testável.
%  A contiguidade hipótese–método é essencial para a
%  coerência dedutiva exigida em periódicos Q1.
% ============================================================

Diante dessa complexidade estrutural, formula-se a questão central desta pesquisa: como evolui espacialmente a distribuição de vento e precipitação nos ciclones extratropicais do Atlântico Sul, e quais configurações caracterizam essa organização durante as fases incipiente, intensificação, maturidade e decaimento? Postula-se que os campos de vento e precipitação apresentam padrões de assimetria sistemática que variam entre fases evolutivas e entre regiões de gênese, e que tais padrões não são capturáveis mediante métricas escalares de intensidade central.

\label{subsec:intro_estrategia_estadistica}
Para testar essa hipótese, são necessárias técnicas estatísticas que transcendam a análise univariada e capturem a covariabilidade entre campos físicos distintos.
Dado que a estrutura ciclônica apresenta uma variabilidade espacial contínua na qual a circulação do vento e os núcleos de precipitação podem covariar durante a evolução do sistema, torna-se necessário recorrer a ferramentas de redução dimensional que preservem a interpretabilidade física. Nesse sentido, \citet{hannachi2007empirical} estabelece a análise de Funções Ortogonais Empíricas Multivariadas (MEOF) como o marco apropriado; ao contrário da análise univariada de EOF aplicada separadamente a cada campo, a técnica MEOF permite extrair modos de variabilidade inerentemente acoplados entre variáveis fisicamente distintas \citep{liang2018multivariate}, capturando assim a organização espacial conjunta de ambos os campos.

A combinação de dados (ERA5) com detecção dinâmica ($\zeta_{850}$), classificação objetiva de fases (\textit{CycloPhaser}) e análise multivariada (MEOF) constitui a base técnica para quantificar a evolução da estrutura espacial dos campos de vento e precipitação na região.
A aplicação desse marco sobre as fases previamente classificadas permitirá construir protótipos estruturais, sintetizando a morfologia típica do ciclone para cada região de gênese e fase evolutiva.

% ============================================================
%  BLOCO F — OBJETIVOS
%  Propósito: Traduzir a pergunta e a hipótese do Bloco E
%  em metas operacionais concretas e verificáveis, com
%  correspondência direta às seções metodológicas
%  do manuscrito.
% ============================================================

\section{Objetivos do Projeto de Pesquisa}

\subsection{Objetivo Geral}
Caracterizar a evolução da estrutura espacial do vento a 10 metros e da precipitação durante as fases do ciclo de vida de ciclones extratropicais no Atlântico Sul, caracterizando seus padrões de variabilidade e a estimação de densidade probabilística de máximos, com estratificação por região de gênese.

\subsection{Objetivos Específicos}
\begin{enumerate}
    \item Consolidar uma base de dados lagrangiana que associe os campos horários de alta resolução de vento e precipitação (ERA5) às trajetórias e fases dinâmicas (incipiente, intensificação, maturidade e decaimento) de ciclones extratropicais do Atlântico Sul, estratificando a amostra por região de gênese (SBR, LPB e ARG) e por intensidade dinâmica.
    
    \item Quantificar a distribuição estatística e a assimetria espacial dos máximos de vento a 10 m e precipitação mediante estimação de densidade por núcleos, caracterizando tanto a distribuição de probabilidade de suas magnitudes (PDF unidimensional) quanto a localização preferencial desses máximos em relação ao centro do ciclone (PDF espacial).
    
    \item Identificar os modos dominantes de variabilidade espacial dos campos de vento e precipitação mediante Funções Ortogonais Empíricas univariadas (EOF) e multivariadas (MEOF), avaliando o padrão estrutural individual de cada variável e a covariabilidade acoplada.
    
    \item Construir protótipos estruturais sintéticos por fase evolutiva e região de gênese, integrando os campos médios condicionados ao modo dominante de covariabilidade (MEOF) com as densidades probabilísticas de localização de máximos (PDFe), para quantificar a magnitude característica de vento e precipitação.
\end{enumerate}